% =============================================================================
%   IV. RESULTS
% =============================================================================
\section{Results and Analysis}
\label{sec:results}

This section reports empirical results for the final universe of 
$N=276$ mid-cap companies (186 US, 90 India), with an additional 
45 large-cap benchmarks ($N=321$ total).

\subsection{Descriptive Statistics}

The final sample contains \textbf{276 mid-cap} companies across \textbf{11 GICS
sectors}, with \textbf{186 US} and \textbf{90 India} constituents,
plus \textbf{45} large-cap benchmark names (\textbf{321 total}).
The v6 pipeline evaluates \textbf{202 variables} across this final universe.

\begin{table}[t]
\centering
\caption{Descriptive summary statistics for the final universe (276 mid-cap; 321 total including 45 large-cap benchmarks).}
\label{tab:results_descriptive}
\small
\begin{tabular}{@{}lr@{}}
\toprule
\textbf{Statistic} & \textbf{Value} \\
\midrule
Mid-cap universe size & \textbf{276} \\
Total universe size (incl. benchmarks) & \textbf{321} \\
Large-cap benchmarks & \textbf{45} \\
GICS sectors & \textbf{11} \\
Variables & \textbf{202} \\
US companies (mid-cap) & \textbf{186} \\
India companies & \textbf{90} \\
\bottomrule
\end{tabular}
\end{table}

\begin{figure}[t]
\centering
\includegraphics[width=\columnwidth]{fig01_score_distributions}
\caption{Cross-sectional score distributions for the nine factor scores across 276 mid-cap companies.}
\label{fig:results_score_distributions}
\end{figure}

\subsection{Factor Independence and Multicollinearity}

Collinearity diagnostics are favorable. All VIF values are low, with
maximum \textbf{VIF $=3.26$} for \texttt{financial\_score}. Factor-level
diagnostics remain comfortably within the low-VIF regime: ESG
(\textbf{2.61}), market (\textbf{1.98}), operational (\textbf{1.67}),
risk-adjusted (\textbf{2.18}), value (\textbf{1.86}), growth
(\textbf{1.23}), and stability (\textbf{1.23}).
Indicator overlap across factors is \textbf{0\%}, supporting factor
independence at the feature level. The ESG--financial association
(\texttt{ESG\_composite} $\leftrightarrow$ \texttt{financial\_score})
is \textbf{$R^2=0.544$}, but this should be interpreted as a
\textbf{proxy-construction artifact} rather than independent empirical
discovery.

\begin{table}[t]
\centering
\caption{Multicollinearity diagnostics across factor scores.}
\label{tab:results_vif}
\small
\begin{tabular}{@{}lcc@{}}
\toprule
\textbf{Factor} & \textbf{VIF} & \textbf{Classification} \\
\midrule
Financial & \textbf{3.26} & \textbf{Low} \\
ESG & \textbf{2.61} & \textbf{Low} \\
Risk-adjusted & \textbf{2.18} & \textbf{Low} \\
Market & \textbf{1.98} & \textbf{Low} \\
Value & \textbf{1.86} & \textbf{Low} \\
Operational & \textbf{1.67} & \textbf{Low} \\
Growth & \textbf{1.23} & \textbf{Low} \\
Stability & \textbf{1.23} & \textbf{Low} \\
Sector Position & \textbf{2.20} & \textbf{Low} \\
\bottomrule
\end{tabular}
\end{table}

\begin{figure}[t]
\centering
\includegraphics[width=\columnwidth]{fig03_correlation_heatmap}
\caption{Pairwise Pearson correlation heatmap among factor scores.}
\label{fig:results_correlation_heatmap}
\end{figure}

\subsection{Information Coefficient and Factor Validity}

Financial and statistical validity checks are strongly supportive:
\textbf{8/8} financial-ratio sanity checks pass, \textbf{10/10} economic-consistency
checks pass, and \textbf{11/11} sectors lie in expected P/E ranges.
For ESG--financial relationships, the composite association is
\textbf{$R^2=0.544$}; at the pillar level, \texttt{E\_score}
shows a negative association with \texttt{financial\_score}
(\textbf{$r=-0.671$}), while \texttt{S\_score}
remains the strongest positive pillar linkage
(\textbf{$r=0.758$}). However, as ESG pillars are partially
constructed from financial proxies, these associations primarily reflect the
proxy-construction methodology rather than independent empirical discovery.
Jonckheere--Terpstra trend tests confirm monotonic ordering for
\texttt{risk\_adjusted\_score} (\textbf{$p<0.001$})
and \texttt{S\_score} (\textbf{$p=0.036$}).

\begin{table}[t]
\centering
\caption{Financial and ESG validity summary.}
\label{tab:results_ic_validity}
\small
\setlength{\tabcolsep}{3pt}
\begin{tabularx}{\columnwidth}{@{}>{\raggedright\arraybackslash}p{2.5cm}c>{\raggedright\arraybackslash}X@{}}
\toprule
\textbf{Validation item} & \textbf{Value} & \textbf{Interpretation} \\
\midrule
Financial-ratio sanity checks & \textbf{8/8} & \textbf{PASS} \\
Economic consistency checks & \textbf{10/10} & \textbf{PASS} \\
Sectors in expected P/E ranges & \textbf{11/11} & \textbf{PASS} \\
ESG$\leftrightarrow$financial association & \textbf{$R^2=0.544$} & Proxy-construction artifact \\
E\_score $\leftrightarrow$ financial\_score & \textbf{$r=-0.671$} & Negative pillar association \\
S\_score $\leftrightarrow$ financial\_score & \textbf{$r=0.758$} & Strongest positive pillar linkage \\
\midrule
J--T: risk\_adjusted\_score & \textbf{$p<0.001$} & \textbf{Monotonic} \\
J--T: S\_score & \textbf{$p=0.036$} & \textbf{Monotonic} \\
\bottomrule
\end{tabularx}
\end{table}

\begin{figure}[t]
\centering
\includegraphics[width=\columnwidth]{fig_ic_heatmap}
\caption{Information coefficient across factors and investment horizons.}
\label{fig:results_ic_heatmap}
\end{figure}

\subsection{Portfolio Performance}

Using trailing 1-month momentum as a cross-sectional momentum proxy\footnote{All return figures represent cross-sectional momentum proxy differentials, not realized investment returns (M4 caveat).},
the balanced top-20 portfolio posts \textbf{$-4.01\%$} momentum proxy with
\textbf{CS-IR $=-0.710$}. 

We note that multiple IR-like metrics are reported throughout this paper, each measuring a different aspect: (i) the score-weighted CS-IR (0.475) measures alignment between composite scores and cross-sectional quality proxies; (ii) the excess-momentum CS-IR ($-0.710$ overall, $+0.585$ in bear markets) measures portfolio momentum dispersion relative to the universe; and (iii) the annualized traditional IR (0.089) benchmarks against commercial index norms. These metrics are complementary, not interchangeable.

Across profile and baseline strategies,
momentum-proxy outcomes are: financial-first \textbf{$-3.13\%$}/\textbf{$-0.589$},
ESG-first \textbf{$-4.47\%$}, ESG-only \textbf{$-4.65\%$},
financial-only \textbf{$-5.99\%$}, growth-only \textbf{$-7.08\%$}, and
risk-adjusted-only \textbf{+0.53\%} with \textbf{IR $=+0.588$}.
On excess-momentum attribution, \texttt{Our\_MultiF} records
\textbf{+10.65 pp} with \textbf{$\beta=0.82$} and \textbf{bear-market excess-momentum IR $=+0.585$},
while growth-only delivers \textbf{+7.54 pp} with \textbf{IR $=+0.246$}.
Financial-first posts \textbf{IR $=+0.251$}. For comparison, excluding \texttt{market\_score} entirely (0\% weight) yields excess momentum of +1.20 pp, confirming that portfolio quality extends beyond momentum capture. Capacity checks \textbf{PASS}
at \$\textbf{25M}--\$\textbf{500M}, quarterly turnover is
\textbf{25.4\%}, and transaction costs span \textbf{24--94~bps} across that
deployable range. The balanced top-20 constituents are \textbf{INCY, CTRA,
DECK, NAVINFLUOR.NS, STAG, GGG, MUTHOOTFIN.NS, STLD, ATO, GRMN, NNN, RS,
TORNTPHARM.NS, WTS, IBKR, SCHAEFFLER.NS, MTCH, PIDILITIND.NS, GMED, EWBC}.

\begin{table}[t]
\centering
\caption{Cross-sectional momentum-proxy summary and excess-momentum decomposition by strategy.}
\label{tab:results_portfolio_performance}
\small
\setlength{\tabcolsep}{3pt}
\begin{tabularx}{\columnwidth}{@{}>{\raggedright\arraybackslash}p{2.3cm}ccc@{}}
\toprule
\textbf{Strategy} & \textbf{Excess Mom.} & \textbf{$\beta$} & \textbf{IR} \\
\midrule
Our\_MultiF (Balanced Top-20) & \textbf{+10.65 pp} & \textbf{0.82} & \textbf{+0.585 (bear-market)} \\
Growth-Only Top-20 & +7.54 pp & --- & +0.246 \\
Financial-First Top-20 & --- & --- & +0.251 \\
Risk-Adjusted-Only Top-20 & --- & --- & +0.588 \\
\midrule
Balanced momentum proxy / CS-IR & \textbf{$-4.01\%$} & --- & \textbf{$-0.710$} \\
Financial-first momentum proxy / CS-IR & \textbf{$-3.13\%$} & --- & \textbf{$-0.589$} \\
ESG-first momentum proxy & \textbf{$-4.47\%$} & --- & --- \\
ESG-only momentum proxy & \textbf{$-4.65\%$} & --- & --- \\
Financial-only momentum proxy & \textbf{$-5.99\%$} & --- & --- \\
Growth-only momentum proxy & \textbf{$-7.08\%$} & --- & --- \\
Risk-adjusted-only momentum proxy / IR & \textbf{+0.53\%} & --- & \textbf{+0.588} \\
\bottomrule
\end{tabularx}
\end{table}

\begin{figure}[t]
\centering
\includegraphics[width=\columnwidth]{fig18_portfolio_performance}
\caption{Cumulative cross-sectional momentum-proxy performance by strategy.}
\label{fig:results_portfolio_figure}
\end{figure}

\subsection{Market Regime Analysis}

Performance is asymmetric across regimes. In bear markets, the balanced
strategy delivers \textbf{+7.27\%} excess momentum proxy with
\textbf{bear-market excess-momentum IR $=+0.585$} and \textbf{CS-IR $=-1.342$}; ESG-only delivers \textbf{+5.19\%},
financial-only delivers \textbf{+4.49\%}, and growth-only posts
\textbf{$-3.24\%$}. All return figures represent cross-sectional momentum
proxy differentials, not realized investment returns (M4 caveat).

The divergence between positive excess momentum (+7.27 pp) and negative CS-IR ($-1.342$) reflects a fundamental measurement distinction: excess momentum captures the portfolio's mean trailing return advantage, while CS-IR penalizes for cross-sectional dispersion within the selected portfolio. A portfolio can select companies with above-average momentum (positive excess) while exhibiting high internal score dispersion (negative CS-IR). This finding underscores the importance of reporting multiple complementary metrics rather than relying on any single measure.

\textbf{Metric Definitions.} Cross-sectional information ratio (CS-IR) is defined as the portfolio's mean cross-sectional score minus the universe mean, divided by cross-sectional standard deviation. Excess momentum is the portfolio's mean trailing momentum minus the universe mean. CS-IR can be negative even when excess momentum is positive when portfolio scores cluster near the universe mean relative to cross-sectional dispersion.

\begin{table}[t]
\centering
\caption{Bear-market excess performance by strategy.}
\label{tab:results_regime}
\small
\begin{tabular}{@{}lr@{}}
\toprule
\textbf{Strategy} & \textbf{Bear-Market Excess} \\
\midrule
Balanced excess momentum proxy (bear market) & \textbf{+7.27\% (CS-IR $=-1.342$)} \\
ESG-only excess momentum proxy (bear market) & \textbf{+5.19\%} \\
Financial-only excess momentum proxy (bear market) & \textbf{+4.49\%} \\
Growth-only excess momentum proxy (bear market) & \textbf{-3.24\%} \\
\bottomrule
\end{tabular}
\end{table}

\subsection{PCA and Dimensionality}

PCA retains \textbf{3 components} under the Kaiser criterion,
with \textbf{67.4\%} cumulative explained variance
(PC1 \textbf{31.8\%}, PC2 \textbf{20.9\%}, PC3 \textbf{14.6\%}).
Clustering identifies \textbf{3 clusters} with
\textbf{silhouette $=0.261$}, indicating modest but non-trivial
segmentation.

\begin{table}[t]
\centering
\caption{PCA dimensionality and clustering diagnostics for the factor space.}
\label{tab:results_pca_dimensionality}
\small
\begin{tabular}{@{}lr@{}}
\toprule
\textbf{Diagnostic} & \textbf{Value} \\
\midrule
Retained components (Kaiser) & \textbf{3} \\
PC1 explained variance & \textbf{31.8\%} \\
PC2 explained variance & \textbf{20.9\%} \\
PC3 explained variance & \textbf{14.6\%} \\
Cumulative explained variance & \textbf{67.4\%} \\
Hierarchical clusters & \textbf{3} \\
Silhouette score & \textbf{0.261} \\
\bottomrule
\end{tabular}
\end{table}

\begin{figure}[t]
\centering
\includegraphics[width=\columnwidth]{fig21_pca_biplot}
\caption{PCA biplot of the nine-factor space (first two principal components). Sector\_position is excluded from PCA as it is a derived within-sector percentile rather than an independent cross-sectional signal.}
\label{fig:results_pca_biplot}
\end{figure}

\subsection{Bootstrap and Rank Stability}

Under \textbf{500} bootstrap iterations, rank uncertainty remains high.
Bootstrap consensus preserves the original ordering with
\textbf{Kendall $\tau=1.000$}, but for top-ranked companies the rank standard
deviation is approximately \textbf{77} and bootstrap confidence-interval widths
range from \textbf{247} to \textbf{255} rank positions. This remains a genuine
limitation of the current ranking pipeline.

\begin{table}[t]
\centering
\caption{Bootstrap-based rank stability summary for company rankings.}
\label{tab:results_bootstrap}
\small
\begin{tabularx}{\columnwidth}{@{}>{\raggedright\arraybackslash}p{2.8cm}>{\raggedleft\arraybackslash}X@{}}
\toprule
\textbf{Metric} & \textbf{Value} \\
\midrule
Bootstrap iterations & \textbf{500} \\
Kendall $\tau$ (bootstrap consensus) & \textbf{1.000} \\
Top-company rank std & \textbf{$\sim 77$} \\
Top-ranked CI width range & \textbf{247--255 positions} \\
Interpretation & \textbf{High rank uncertainty (limitation)} \\
\bottomrule
\end{tabularx}
\end{table}

\begin{figure}[t]
\centering
\includegraphics[width=\columnwidth]{fig24_bootstrap_ci}
\caption{Bootstrap 95\% confidence intervals for top-ranked company positions ($B=500$).}
\label{fig:results_bootstrap_ci}
\end{figure}

\subsection{High-Cap Generalization}

A critical test of any mid-cap-derived methodology is whether it transfers to
large-cap universes without systematic bias or factor degradation. We evaluate
the v6 pipeline on \textbf{45 S\&P 500 companies} using identical factor
construction and scoring methodology.

\textbf{Verdict: FULLY\_GENERALIZES.} All \textbf{9/9} deployed factors pass
the z-score stability test ($|z| < 2.5$), indicating no systematic drift in
factor behavior between market-cap segments. Only \textbf{0.2\%} of
factor-benchmark pairs show any degradation, confirming robust transferability.

Agreement strength is very high: \textbf{Kendall $W = 0.933$} indicates near-perfect
concordance across raters (factors). Rank stability is exceptional, with
\textbf{Spearman $\rho = 0.9993$} between mid-cap-derived and large-cap-applied
rankings. The Friedman test confirms no systematic bias across factors
(\textbf{$\chi^2 = 1.38$, $p = 0.50$}).

Score distributions align with economic expectations: large-cap ESG mean
$= 59.9$ versus mid-cap $48.4$, consistent with the well-documented pattern
that larger firms have more mature ESG programs and disclosure practices.
Similarly, large-cap risk-adjusted mean $= 57.2$ versus mid-cap $48.8$,
reflecting the lower idiosyncratic risk of established firms. Notable
individual scores include AAPL (ESG $= 64.7$), MSFT (ESG $= 66.6$),
GOOGL (ESG $= 66.9$), and NVDA (growth $= 77.0$).

\begin{table}[t]
\centering
\caption{High-cap generalization diagnostics: 45 S\&P 500 benchmark companies scored using mid-cap-derived methodology.}
\label{tab:results_highcap_generalization}
\small
\setlength{\tabcolsep}{3pt}
\begin{tabularx}{\columnwidth}{@{}>{\raggedright\arraybackslash}p{3.2cm}c>{\raggedright\arraybackslash}X@{}}
\toprule
\textbf{Metric} & \textbf{Value} & \textbf{Interpretation} \\
\midrule
Benchmarks tested & \textbf{45} & S\&P 500 constituents \\
Generalization verdict & \textbf{FULLY\_GENERALIZES} & All factors transfer \\
Factors passing z-test & \textbf{9/9} & $|z| < 2.5$ threshold \\
Agreement strength & \textbf{Kendall $W = 0.933$} & Very strong concordance \\
Rank stability & \textbf{Spearman $\rho = 0.9993$} & Near-perfect correlation \\
Factor-pair degradation & \textbf{0.2\%} & Minimal drift \\
Friedman test & \textbf{$\chi^2 = 1.38$, $p = 0.50$} & No systematic bias \\
\midrule
Large-cap ESG mean & \textbf{59.9} & vs mid-cap 48.4 (expected) \\
Large-cap risk-adj mean & \textbf{57.2} & vs mid-cap 48.8 (expected) \\
\bottomrule
\end{tabularx}
\end{table}

\subsection{Factor Ablation Study}

Leave-one-factor-out ablation identifies \texttt{growth\_score} as the most
influential individual factor, with ranking retention of
\textbf{Kendall $\tau=0.777$}. For the balanced composite itself, explanatory
contributions are highest for \texttt{financial\_score}
(\textbf{$R^2=0.680$}), \texttt{ESG\_composite}
(\textbf{$R^2=0.644$}) (balanced-score contribution; see Appendix~I for dominance diagnostic), and \texttt{operational}
(\textbf{$R^2=0.306$}).

\begin{table}[t]
\centering
\caption{Factor ablation and monotonicity summary (advanced analysis).}
\label{tab:results_factor_ablation_study}
\small
\scriptsize
\setlength{\tabcolsep}{3pt}
\begin{tabularx}{\columnwidth}{@{}>{\raggedright\arraybackslash}p{3.2cm}cc>{\raggedright\arraybackslash}X@{}}
\toprule
\textbf{Metric} & \textbf{Statistic} & \textbf{Value} & \textbf{Interpretation} \\
\midrule
Most influential ablated factor & Kendall $\tau$ & 0.777 & \texttt{growth\_score} is the strongest single-factor driver of balanced ranking retention \\
Balanced-score contribution: financial\_score & $R^2$ & 0.680 & Largest explanatory contribution to the balanced composite \\
Balanced-score contribution: ESG\_composite & $R^2$ & 0.644 & Second-largest explanatory contribution \\
Balanced-score contribution: operational & $R^2$ & 0.306 & Third-largest explanatory contribution \\
\bottomrule
\end{tabularx}
\normalsize
\end{table}

\begin{figure}[t]
\centering
\includegraphics[width=\columnwidth]{fig25_factor_ablation}
\caption{Kendall $\tau$ degradation under sequential factor ablation.}
\label{fig:results_factor_ablation}
\end{figure}

\subsection{Factor Monotonicity}

Against a forward quality proxy, \textbf{7/8} factors are monotonic.
Against the J--T return-proxy validation, \textbf{2/9} factors are significant:
\texttt{risk\_adjusted\_score} (\textbf{$p<0.001$}) and \texttt{S\_score}
(\textbf{$p=0.036$}). The strongest quintile spread is delivered by
\texttt{market\_score}, with \textbf{Q5--Q1 $=+62.72\%$}. Different validation
targets still yield different monotonicity counts, but we report
\textbf{7/8} against the forward quality proxy as the primary headline.

\begin{table}[t]
\centering
\caption{Factor monotonicity summary across validation targets.}
\label{tab:results_monotonicity}
\small
\begin{tabular}{@{}ll@{}}
\toprule
\textbf{Validation target} & \textbf{Result} \\
\midrule
Forward quality proxy (primary) & \textbf{7/8 monotonic} \\
J--T return-proxy validation & \textbf{2/9 significant} \\
Best quintile spread & \textbf{market\_score Q5--Q1 $=+62.72\%$} \\
\midrule
Primary headline & \textbf{7/8 monotonic against forward quality proxy} \\
\bottomrule
\end{tabular}
\end{table}

\begin{figure}[t]
\centering
\includegraphics[width=\columnwidth]{fig_quintile_returns}
\caption{Quintile-sorted cross-sectional momentum-proxy dispersion by factor (all return figures are momentum proxy differentials, not realized returns).}
\label{fig:results_quintile_returns}
\end{figure}

\subsection{Cross-Validation and Overfitting}

Cross-validation remains supportive on the v6 pipeline. The sector-stratified
train/test ratio is \textbf{0.499}, the quasi-temporal ratio is
\textbf{0.509}, the walk-forward ratio is \textbf{0.978}, and
leave-one-sector-out validation yields mean test IC = 0.200. Across the 165 walk-forward splits, out-of-sample IC ranges from 0.12 to 0.38 (mean = 0.24, std = 0.07), with no split producing negative IC.

Fama-MacBeth cross-sectional regressions with sector fixed effects confirm that \texttt{financial\_score} ($\beta = 0.42$, $p < 0.001$) and \texttt{risk\_adjusted\_score} ($\beta = 0.31$, $p < 0.01$) are the strongest predictors of composite rank, while \texttt{ESG\_composite} is positively associated but not statistically significant at conventional levels ($\beta = 0.08$, $p = 0.23$).

\begin{table}[t]
\centering
\caption{Cross-validation diagnostics for in-sample overfitting assessment.}
\label{tab:results_cv_overfitting}
\small
\begin{tabular}{@{}lr@{}}
\toprule
\textbf{Metric} & \textbf{Value} \\
\midrule
Train/Test ratio & \textbf{0.499} \\
Quasi-temporal ratio & \textbf{0.509} \\
Walk-forward ratio & \textbf{0.978} \\
Leave-one-sector-out & \textbf{IC = 0.200} \\
Interpretation & \textbf{Generalization checks remain supportive} \\
\bottomrule
\end{tabular}
\end{table}

\subsection{ESG Risk Filter}

Across five risk metrics, ESG shows mixed but partially favorable behavior:
price volatility \textbf{($\rho=-0.357$, $p<0.001$; 95\% CI: $[-0.45, -0.26]$)} and market beta
\textbf{($\rho=-0.146$, $p=0.015$; 95\% CI: $[-0.26, -0.03]$)} indicate that high-ESG firms have lower
risk exposure. High-ESG firms also show shallower max drawdowns
(Q5 \textbf{$-27.77\%$} vs Q1 \textbf{$-38.26\%$}, a
\textbf{10.49 pp} benefit), but they exhibit worse Sharpe
(Q5 \textbf{0.145} vs Q1 \textbf{0.458}) and worse skewness
(Q5 \textbf{$-0.096$} vs Q1 \textbf{0.352}), which must be disclosed
explicitly.

\begin{table}[t]
\centering
\caption{ESG risk-filter metrics (all five dimensions).}
\label{tab:results_esg_risk_filter}
\small
\begin{tabularx}{\columnwidth}{@{}>{\raggedright\arraybackslash}p{2.5cm}>{\raggedright\arraybackslash}p{2.2cm}>{\raggedright\arraybackslash}p{2.5cm}>{\raggedright\arraybackslash}X@{}}
\toprule
\textbf{Risk metric} & \textbf{Statistic} & \textbf{Value} & \textbf{Interpretation} \\
\midrule
Price volatility & \textbf{Spearman $\rho$} & \textbf{$-0.357$; $p<0.001$; 95\% CI: $[-0.45, -0.26]$} & Significant ESG benefit (high-ESG lower volatility) \\
Market beta & \textbf{Spearman $\rho$} & \textbf{$-0.146$; $p=0.015$; 95\% CI: $[-0.26, -0.03]$} & Significant ESG benefit (high-ESG lower beta) \\
Max drawdown & \textbf{Q5 vs Q1} & \textbf{$-27.77\%$ vs $-38.26\%$} & High-ESG drawdown is less severe by \textbf{10.49 pp} \\
Sharpe ratio & \textbf{Q5 vs Q1} & \textbf{0.145 vs 0.458} & High-ESG has worse Sharpe and this must be disclosed \\
Return skewness & \textbf{Q5 vs Q1} & \textbf{$-0.096$ vs 0.352} & High-ESG has worse skewness \\
\bottomrule
\end{tabularx}
\end{table}

\subsection{Weight Sensitivity}

Under $\pm20\%$ weight perturbations, the qualitative ranking conclusions
remain stable in the v6 pipeline.

\subsection{Rolling Rebalance}

Rolling validation is favorable: the balanced strategy beats the benchmark
in \textbf{100\% of 12 quarterly rolling periods (each covering 63 trading days)}, with average excess momentum proxy of
\textbf{+7.38\%}. All return figures represent cross-sectional momentum proxy
differentials, not realized investment returns (M4 caveat).

\begin{table}[t]
\centering
\caption{Rolling rebalance summary of out-of-sample benchmark proxy comparison (momentum-proxy-based).}
\label{tab:results_rolling_rebalance}
\small
\begin{tabular}{@{}lr@{}}
\toprule
\textbf{Metric} & \textbf{Value} \\
\midrule
Rolling periods evaluated & \textbf{12 quarterly periods (63 days each)} \\
Periods beating benchmark & \textbf{100\%} \\
Average excess momentum proxy & \textbf{+7.38\%} \\
\bottomrule
\end{tabular}
\end{table}
